
\documentclass{article}
\usepackage{colortbl}
\usepackage{makecell}
\usepackage{multirow}
\usepackage{supertabular}

\begin{document}

\newcounter{utterance}

\centering \large Interaction Transcript for game `matchit\_ascii', experiment `different\_grid', episode 1 with claude{-}opus{-}4{-}8{-}azure.
\vspace{24pt}

{ \footnotesize  \setcounter{utterance}{1}
\setlength{\tabcolsep}{0pt}
\begin{supertabular}{c@{$\;$}|p{.15\linewidth}@{}p{.15\linewidth}p{.15\linewidth}p{.15\linewidth}p{.15\linewidth}p{.15\linewidth}}
    \# & $\;$A & \multicolumn{4}{c}{Game Master} & $\;\:$B\\
    \hline

    \theutterance \stepcounter{utterance}  
    & & \multicolumn{4}{p{0.6\linewidth}}{
        \cellcolor[rgb]{0.9,0.9,0.9}{
            \makecell[{{p{\linewidth}}}]{
                \texttt{\tiny{[P1$\langle$GM]}}
                \texttt{Zúčastňujete sa spoločnej hry na hádanie.} \\
\texttt{Cieľom je zistiť, či je táto mriežka a iný obrázok, ktorý vidím len ja, rovnaké.} \\
\texttt{Najskôr opíšte svoju mriežku. Potom poskytnem svoj opis a môžeme si navzájom klásť otázky o mriežkach, aby sme zistili, či sú rovnaké. Teraz začnite svoj krátky opis mriežky s "POPIS:", za ktorým nasleduje opis. Nepridávajte nič iné. Toto je vaša mriežka:} \\
\texttt{X X X X X} \\
\texttt{▢ ▢ ▢ ▢ X} \\
\texttt{▢ ▢ X X X} \\
\texttt{▢ ▢ ▢ ▢ X} \\
\texttt{X X X X X} \\
            }
        }
    }
    & \\ \\

    \theutterance \stepcounter{utterance}  
    & \multicolumn{4}{p{0.6\linewidth}}{
        \cellcolor[rgb]{0.9,0.9,0.9}{
            \makecell[{{p{\linewidth}}}]{
                \texttt{\tiny{[P1$\rangle$GM]}}
                \texttt{POPIS: Mriežka má rozmery 5×5. Prvý riadok je celý vyplnený (X X X X X). Druhý riadok má štyri prázdne políčka a posledné X (▢ ▢ ▢ ▢ X). Tretí riadok má dve prázdne políčka a tri X (▢ ▢ X X X). Štvrtý riadok je rovnaký ako druhý (▢ ▢ ▢ ▢ X). Piaty riadok je celý vyplnený (X X X X X). Vyplnené políčka tvoria tvar podobný písmenu, s plným pravým stĺpcom a plnými horným aj dolným riadkom.} \\
            }
        }
    }
    & & \\ \\

    \theutterance \stepcounter{utterance}  
    & & & \multicolumn{2}{p{0.3\linewidth}}{
        \cellcolor[rgb]{0.9,0.9,0.9}{
            \makecell[{{p{\linewidth}}}]{
                \texttt{\tiny{[GM$|$GM]}}
                \texttt{continue} \\
            }
        }
    }
    & & \\ \\

    \theutterance \stepcounter{utterance}  
    & & \multicolumn{4}{p{0.6\linewidth}}{
        \cellcolor[rgb]{0.9,0.9,0.9}{
            \makecell[{{p{\linewidth}}}]{
                \texttt{\tiny{[GM$\rangle$P2]}}
                \texttt{Zúčastňujete sa spoločnej hry na hádanie.} \\
\texttt{Cieľom je zistiť, či je táto mriežka a iný obrázok, ktorý vidím len ja, rovnaké.} \\
\texttt{Najskôr opíšte svoju mriežku. Potom poskytnem svoj opis a môžeme si navzájom klásť otázky o mriežkach, aby sme zistili, či sú rovnaké. Teraz začnite svoj krátky opis mriežky s "POPIS:", za ktorým nasleduje opis. Nepridávajte nič iné. Toto je vaša mriežka:} \\
\texttt{X ▢ ▢ ▢ X} \\
\texttt{X ▢ ▢ X X} \\
\texttt{X ▢ X ▢ X} \\
\texttt{X X ▢ ▢ X} \\
\texttt{X ▢ ▢ ▢ X} \\
            }
        }
    }
    & \\ \\

    \theutterance \stepcounter{utterance}  
    & & & \multicolumn{4}{p{0.6\linewidth}}{
        \cellcolor[rgb]{0.9,0.9,0.9}{
            \makecell[{{p{\linewidth}}}]{
                \texttt{\tiny{[GM$\langle$P2]}}
                \texttt{POPIS: Mriežka 5x5. Prvý aj posledný stĺpec sú úplne vyplnené X. Vnútri tvoria X{-}ká diagonálu: v druhom riadku je X v 4. stĺpci, v treťom riadku v 3. stĺpci, vo štvrtom riadku v 2. stĺpci. Prvý a posledný riadok majú vyplnené len okraje. Ostatné bunky sú prázdne.} \\
            }
        }
    }
     \\ \\

    \theutterance \stepcounter{utterance}  
    & & & \multicolumn{2}{p{0.3\linewidth}}{
        \cellcolor[rgb]{0.9,0.9,0.9}{
            \makecell[{{p{\linewidth}}}]{
                \texttt{\tiny{[GM$|$GM]}}
                \texttt{continue} \\
            }
        }
    }
    & & \\ \\

    \theutterance \stepcounter{utterance}  
    & & \multicolumn{4}{p{0.6\linewidth}}{
        \cellcolor[rgb]{0.9,0.9,0.9}{
            \makecell[{{p{\linewidth}}}]{
                \texttt{\tiny{[GM$\rangle$P2]}}
                \texttt{Toto je môj POPISPOPIS: Mriežka má rozmery 5×5. Prvý riadok je celý vyplnený (X X X X X). Druhý riadok má štyri prázdne políčka a posledné X (▢ ▢ ▢ ▢ X). Tretí riadok má dve prázdne políčka a tri X (▢ ▢ X X X). Štvrtý riadok je rovnaký ako druhý (▢ ▢ ▢ ▢ X). Piaty riadok je celý vyplnený (X X X X X). Vyplnené políčka tvoria tvar podobný písmenu, s plným pravým stĺpcom a plnými horným aj dolným riadkom.} \\
\texttt{Teraz položte otázku, aby ste zistili nové aspekty mojej mriežky, ktoré sa môžu líšiť od vašej mriežky. Začnite s "OTÁZKA:" a nepridávajte nič iné.} \\
            }
        }
    }
    & \\ \\

    \theutterance \stepcounter{utterance}  
    & & & \multicolumn{4}{p{0.6\linewidth}}{
        \cellcolor[rgb]{0.9,0.9,0.9}{
            \makecell[{{p{\linewidth}}}]{
                \texttt{\tiny{[GM$\langle$P2]}}
                \texttt{OTÁZKA: Je vo vašej mriežke celý prvý (ľavý) stĺpec vyplnený X, alebo je v prvom stĺpci len jedno X v prvom a poslednom riadku?} \\
            }
        }
    }
     \\ \\

    \theutterance \stepcounter{utterance}  
    & & & \multicolumn{2}{p{0.3\linewidth}}{
        \cellcolor[rgb]{0.9,0.9,0.9}{
            \makecell[{{p{\linewidth}}}]{
                \texttt{\tiny{[GM$|$GM]}}
                \texttt{continue} \\
            }
        }
    }
    & & \\ \\

    \theutterance \stepcounter{utterance}  
    & & \multicolumn{4}{p{0.6\linewidth}}{
        \cellcolor[rgb]{0.9,0.9,0.9}{
            \makecell[{{p{\linewidth}}}]{
                \texttt{\tiny{[P1$\langle$GM]}}
                \texttt{Toto je môj POPISPOPIS: Mriežka 5x5. Prvý aj posledný stĺpec sú úplne vyplnené X. Vnútri tvoria X{-}ká diagonálu: v druhom riadku je X v 4. stĺpci, v treťom riadku v 3. stĺpci, vo štvrtom riadku v 2. stĺpci. Prvý a posledný riadok majú vyplnené len okraje. Ostatné bunky sú prázdne.} \\
\texttt{OTÁZKA: Je vo vašej mriežke celý prvý (ľavý) stĺpec vyplnený X, alebo je v prvom stĺpci len jedno X v prvom a poslednom riadku?Svoju odpoveď začnite s "ODPOVEĎ:" a nepridávajte nič iné.} \\
            }
        }
    }
    & \\ \\

    \theutterance \stepcounter{utterance}  
    & \multicolumn{4}{p{0.6\linewidth}}{
        \cellcolor[rgb]{0.9,0.9,0.9}{
            \makecell[{{p{\linewidth}}}]{
                \texttt{\tiny{[P1$\rangle$GM]}}
                \texttt{ODPOVEĎ: V mojej mriežke nie je celý ľavý stĺpec vyplnený. V prvom stĺpci sú X iba v prvom a poslednom riadku (row1 a row5), ostatné bunky prvého stĺpca (row2, row3, row4) sú prázdne. Toto sa líši od vášho popisu, kde je celý prvý stĺpec vyplnený X. OTÁZKA: Je vo vašom prvom riadku vyplnený iba prvý a posledný stĺpec (X ▢ ▢ ▢ X), alebo je celý prvý riadok vyplnený (X X X X X)?} \\
            }
        }
    }
    & & \\ \\

    \theutterance \stepcounter{utterance}  
    & & & \multicolumn{2}{p{0.3\linewidth}}{
        \cellcolor[rgb]{0.9,0.9,0.9}{
            \makecell[{{p{\linewidth}}}]{
                \texttt{\tiny{[GM$|$GM]}}
                \texttt{continue} \\
            }
        }
    }
    & & \\ \\

    \theutterance \stepcounter{utterance}  
    & & \multicolumn{4}{p{0.6\linewidth}}{
        \cellcolor[rgb]{0.9,0.9,0.9}{
            \makecell[{{p{\linewidth}}}]{
                \texttt{\tiny{[P1$\langle$GM]}}
                \texttt{Teraz položte otázku, aby ste zistili nové aspekty mojej mriežky, ktoré sa môžu líšiť od vašej mriežky. Začnite s "OTÁZKA:" a nepridávajte nič iné.} \\
            }
        }
    }
    & \\ \\

    \theutterance \stepcounter{utterance}  
    & \multicolumn{4}{p{0.6\linewidth}}{
        \cellcolor[rgb]{0.9,0.9,0.9}{
            \makecell[{{p{\linewidth}}}]{
                \texttt{\tiny{[P1$\rangle$GM]}}
                \texttt{OTÁZKA: Vo vašom treťom riadku – ktoré presné bunky sú vyplnené X? Vo mojej mriežke je tretí riadok ▢ ▢ X X X (prázdne prvé dve bunky, potom tri X v stĺpcoch 3, 4 a 5). Aké je to u vás?} \\
            }
        }
    }
    & & \\ \\

    \theutterance \stepcounter{utterance}  
    & & & \multicolumn{2}{p{0.3\linewidth}}{
        \cellcolor[rgb]{0.9,0.9,0.9}{
            \makecell[{{p{\linewidth}}}]{
                \texttt{\tiny{[GM$|$GM]}}
                \texttt{continue} \\
            }
        }
    }
    & & \\ \\

    \theutterance \stepcounter{utterance}  
    & & \multicolumn{4}{p{0.6\linewidth}}{
        \cellcolor[rgb]{0.9,0.9,0.9}{
            \makecell[{{p{\linewidth}}}]{
                \texttt{\tiny{[GM$\rangle$P2]}}
                \texttt{ODPOVEĎ: V mojej mriežke nie je celý ľavý stĺpec vyplnený. V prvom stĺpci sú X iba v prvom a poslednom riadku (row1 a row5), ostatné bunky prvého stĺpca (row2, row3, row4) sú prázdne. Toto sa líši od vášho popisu, kde je celý prvý stĺpec vyplnený X. OTÁZKA: Je vo vašom prvom riadku vyplnený iba prvý a posledný stĺpec (X ▢ ▢ ▢ X), alebo je celý prvý riadok vyplnený (X X X X X)?} \\
\texttt{OTÁZKA: Vo vašom treťom riadku – ktoré presné bunky sú vyplnené X? Vo mojej mriežke je tretí riadok ▢ ▢ X X X (prázdne prvé dve bunky, potom tri X v stĺpcoch 3, 4 a 5). Aké je to u vás?Svoju odpoveď začnite s "ODPOVEĎ:" a nepridávajte nič iné.} \\
            }
        }
    }
    & \\ \\

    \theutterance \stepcounter{utterance}  
    & & & \multicolumn{4}{p{0.6\linewidth}}{
        \cellcolor[rgb]{0.9,0.9,0.9}{
            \makecell[{{p{\linewidth}}}]{
                \texttt{\tiny{[GM$\langle$P2]}}
                \texttt{ODPOVEĎ: V mojom prvom riadku je vyplnený iba prvý a posledný stĺpec (X ▢ ▢ ▢ X), nie celý riadok. V mojom treťom riadku sú X v stĺpcoch 1, 3 a 5 (X ▢ X ▢ X). Naše mriežky sa teda líšia – tieto rozdiely potvrdzujú, že mriežky nie sú rovnaké.} \\
            }
        }
    }
     \\ \\

    \theutterance \stepcounter{utterance}  
    & & & \multicolumn{2}{p{0.3\linewidth}}{
        \cellcolor[rgb]{0.9,0.9,0.9}{
            \makecell[{{p{\linewidth}}}]{
                \texttt{\tiny{[GM$|$GM]}}
                \texttt{continue} \\
            }
        }
    }
    & & \\ \\

    \theutterance \stepcounter{utterance}  
    & & \multicolumn{4}{p{0.6\linewidth}}{
        \cellcolor[rgb]{0.9,0.9,0.9}{
            \makecell[{{p{\linewidth}}}]{
                \texttt{\tiny{[GM$\rangle$P2]}}
                \texttt{Teraz položte otázku, aby ste zistili nové aspekty mojej mriežky, ktoré sa môžu líšiť od vašej mriežky. Začnite s "OTÁZKA:" a nepridávajte nič iné.} \\
            }
        }
    }
    & \\ \\

    \theutterance \stepcounter{utterance}  
    & & & \multicolumn{4}{p{0.6\linewidth}}{
        \cellcolor[rgb]{0.9,0.9,0.9}{
            \makecell[{{p{\linewidth}}}]{
                \texttt{\tiny{[GM$\langle$P2]}}
                \texttt{OTÁZKA: Vo vašom druhom riadku – je vyplnená iba posledná bunka (▢ ▢ ▢ ▢ X), alebo sú vyplnené aj nejaké iné bunky?} \\
            }
        }
    }
     \\ \\

    \theutterance \stepcounter{utterance}  
    & & & \multicolumn{2}{p{0.3\linewidth}}{
        \cellcolor[rgb]{0.9,0.9,0.9}{
            \makecell[{{p{\linewidth}}}]{
                \texttt{\tiny{[GM$|$GM]}}
                \texttt{continue} \\
            }
        }
    }
    & & \\ \\

    \theutterance \stepcounter{utterance}  
    & & \multicolumn{4}{p{0.6\linewidth}}{
        \cellcolor[rgb]{0.9,0.9,0.9}{
            \makecell[{{p{\linewidth}}}]{
                \texttt{\tiny{[P1$\langle$GM]}}
                \texttt{ODPOVEĎ: V mojom prvom riadku je vyplnený iba prvý a posledný stĺpec (X ▢ ▢ ▢ X), nie celý riadok. V mojom treťom riadku sú X v stĺpcoch 1, 3 a 5 (X ▢ X ▢ X). Naše mriežky sa teda líšia – tieto rozdiely potvrdzujú, že mriežky nie sú rovnaké.} \\
\texttt{OTÁZKA: Vo vašom druhom riadku – je vyplnená iba posledná bunka (▢ ▢ ▢ ▢ X), alebo sú vyplnené aj nejaké iné bunky?Svoju odpoveď začnite s "ODPOVEĎ:" a nepridávajte nič iné.} \\
            }
        }
    }
    & \\ \\

    \theutterance \stepcounter{utterance}  
    & \multicolumn{4}{p{0.6\linewidth}}{
        \cellcolor[rgb]{0.9,0.9,0.9}{
            \makecell[{{p{\linewidth}}}]{
                \texttt{\tiny{[P1$\rangle$GM]}}
                \texttt{ODPOVEĎ: V mojom druhom riadku je vyplnená iba posledná bunka (▢ ▢ ▢ ▢ X), teda X iba v 5. stĺpci, ostatné štyri bunky sú prázdne. Vo vašom druhom riadku je X v 1., 4. a 5. stĺpci, takže sa opäť líšia. Na základe všetkých doterajších rozdielov (prvý riadok, druhý riadok, tretí riadok, prvý stĺpec) je jasné, že naše mriežky NIE sú rovnaké.} \\
            }
        }
    }
    & & \\ \\

    \theutterance \stepcounter{utterance}  
    & & & \multicolumn{2}{p{0.3\linewidth}}{
        \cellcolor[rgb]{0.9,0.9,0.9}{
            \makecell[{{p{\linewidth}}}]{
                \texttt{\tiny{[GM$|$GM]}}
                \texttt{continue} \\
            }
        }
    }
    & & \\ \\

    \theutterance \stepcounter{utterance}  
    & & \multicolumn{4}{p{0.6\linewidth}}{
        \cellcolor[rgb]{0.9,0.9,0.9}{
            \makecell[{{p{\linewidth}}}]{
                \texttt{\tiny{[P1$\langle$GM]}}
                \texttt{Teraz položte otázku, aby ste zistili nové aspekty mojej mriežky, ktoré sa môžu líšiť od vašej mriežky. Začnite s "OTÁZKA:" a nepridávajte nič iné.} \\
            }
        }
    }
    & \\ \\

    \theutterance \stepcounter{utterance}  
    & \multicolumn{4}{p{0.6\linewidth}}{
        \cellcolor[rgb]{0.9,0.9,0.9}{
            \makecell[{{p{\linewidth}}}]{
                \texttt{\tiny{[P1$\rangle$GM]}}
                \texttt{OTÁZKA: Vo vašom štvrtom riadku – ktoré presné bunky sú vyplnené X? Vo mojej mriežke je štvrtý riadok ▢ ▢ ▢ ▢ X (prázdne prvé štyri bunky a X iba v 5. stĺpci). Aké je to u vás?} \\
            }
        }
    }
    & & \\ \\

    \theutterance \stepcounter{utterance}  
    & & & \multicolumn{2}{p{0.3\linewidth}}{
        \cellcolor[rgb]{0.9,0.9,0.9}{
            \makecell[{{p{\linewidth}}}]{
                \texttt{\tiny{[GM$|$GM]}}
                \texttt{continue} \\
            }
        }
    }
    & & \\ \\

    \theutterance \stepcounter{utterance}  
    & & \multicolumn{4}{p{0.6\linewidth}}{
        \cellcolor[rgb]{0.9,0.9,0.9}{
            \makecell[{{p{\linewidth}}}]{
                \texttt{\tiny{[GM$\rangle$P2]}}
                \texttt{ODPOVEĎ: V mojom druhom riadku je vyplnená iba posledná bunka (▢ ▢ ▢ ▢ X), teda X iba v 5. stĺpci, ostatné štyri bunky sú prázdne. Vo vašom druhom riadku je X v 1., 4. a 5. stĺpci, takže sa opäť líšia. Na základe všetkých doterajších rozdielov (prvý riadok, druhý riadok, tretí riadok, prvý stĺpec) je jasné, že naše mriežky NIE sú rovnaké.} \\
\texttt{OTÁZKA: Vo vašom štvrtom riadku – ktoré presné bunky sú vyplnené X? Vo mojej mriežke je štvrtý riadok ▢ ▢ ▢ ▢ X (prázdne prvé štyri bunky a X iba v 5. stĺpci). Aké je to u vás?Svoju odpoveď začnite s "ODPOVEĎ:" a nepridávajte nič iné.} \\
            }
        }
    }
    & \\ \\

    \theutterance \stepcounter{utterance}  
    & & & \multicolumn{4}{p{0.6\linewidth}}{
        \cellcolor[rgb]{0.9,0.9,0.9}{
            \makecell[{{p{\linewidth}}}]{
                \texttt{\tiny{[GM$\langle$P2]}}
                \texttt{ODPOVEĎ: V mojom štvrtom riadku sú X v stĺpcoch 1, 2 a 5 (X X ▢ ▢ X). Vo vašom štvrtom riadku je X iba v 5. stĺpci, takže sa opäť líšia. Naše mriežky NIE sú rovnaké.} \\
            }
        }
    }
     \\ \\

    \theutterance \stepcounter{utterance}  
    & & & \multicolumn{2}{p{0.3\linewidth}}{
        \cellcolor[rgb]{0.9,0.9,0.9}{
            \makecell[{{p{\linewidth}}}]{
                \texttt{\tiny{[GM$|$GM]}}
                \texttt{continue} \\
            }
        }
    }
    & & \\ \\

    \theutterance \stepcounter{utterance}  
    & & \multicolumn{4}{p{0.6\linewidth}}{
        \cellcolor[rgb]{0.9,0.9,0.9}{
            \makecell[{{p{\linewidth}}}]{
                \texttt{\tiny{[GM$\rangle$P2]}}
                \texttt{Teraz položte otázku, aby ste zistili nové aspekty mojej mriežky, ktoré sa môžu líšiť od vašej mriežky. Začnite s "OTÁZKA:" a nepridávajte nič iné.} \\
            }
        }
    }
    & \\ \\

    \theutterance \stepcounter{utterance}  
    & & & \multicolumn{4}{p{0.6\linewidth}}{
        \cellcolor[rgb]{0.9,0.9,0.9}{
            \makecell[{{p{\linewidth}}}]{
                \texttt{\tiny{[GM$\langle$P2]}}
                \texttt{OTÁZKA: Vo vašom poslednom (piatom) riadku je celý riadok vyplnený (X X X X X), alebo iba prvý a posledný stĺpec?} \\
            }
        }
    }
     \\ \\

    \theutterance \stepcounter{utterance}  
    & & & \multicolumn{2}{p{0.3\linewidth}}{
        \cellcolor[rgb]{0.9,0.9,0.9}{
            \makecell[{{p{\linewidth}}}]{
                \texttt{\tiny{[GM$|$GM]}}
                \texttt{continue} \\
            }
        }
    }
    & & \\ \\

    \theutterance \stepcounter{utterance}  
    & & \multicolumn{4}{p{0.6\linewidth}}{
        \cellcolor[rgb]{0.9,0.9,0.9}{
            \makecell[{{p{\linewidth}}}]{
                \texttt{\tiny{[P1$\langle$GM]}}
                \texttt{ODPOVEĎ: V mojom štvrtom riadku sú X v stĺpcoch 1, 2 a 5 (X X ▢ ▢ X). Vo vašom štvrtom riadku je X iba v 5. stĺpci, takže sa opäť líšia. Naše mriežky NIE sú rovnaké.} \\
\texttt{OTÁZKA: Vo vašom poslednom (piatom) riadku je celý riadok vyplnený (X X X X X), alebo iba prvý a posledný stĺpec?Svoju odpoveď začnite s "ODPOVEĎ:" a nepridávajte nič iné.} \\
            }
        }
    }
    & \\ \\

    \theutterance \stepcounter{utterance}  
    & \multicolumn{4}{p{0.6\linewidth}}{
        \cellcolor[rgb]{0.9,0.9,0.9}{
            \makecell[{{p{\linewidth}}}]{
                \texttt{\tiny{[P1$\rangle$GM]}}
                \texttt{ODPOVEĎ: V mojom piatom riadku je celý riadok vyplnený (X X X X X), teda X vo všetkých piatich stĺpcoch. Vo vašom piatom riadku je vyplnený iba prvý a posledný stĺpec (X ▢ ▢ ▢ X), takže sa líšia. Na základe všetkých doterajších rozdielov je jednoznačné, že naše mriežky NIE sú rovnaké.} \\
            }
        }
    }
    & & \\ \\

    \theutterance \stepcounter{utterance}  
    & & & \multicolumn{2}{p{0.3\linewidth}}{
        \cellcolor[rgb]{0.9,0.9,0.9}{
            \makecell[{{p{\linewidth}}}]{
                \texttt{\tiny{[GM$|$GM]}}
                \texttt{continue} \\
            }
        }
    }
    & & \\ \\

    \theutterance \stepcounter{utterance}  
    & & \multicolumn{4}{p{0.6\linewidth}}{
        \cellcolor[rgb]{0.9,0.9,0.9}{
            \makecell[{{p{\linewidth}}}]{
                \texttt{\tiny{[P1$\langle$GM]}}
                \texttt{Teraz položte otázku, aby ste zistili nové aspekty mojej mriežky, ktoré sa môžu líšiť od vašej mriežky. Začnite s "OTÁZKA:" a nepridávajte nič iné.} \\
            }
        }
    }
    & \\ \\

    \theutterance \stepcounter{utterance}  
    & \multicolumn{4}{p{0.6\linewidth}}{
        \cellcolor[rgb]{0.9,0.9,0.9}{
            \makecell[{{p{\linewidth}}}]{
                \texttt{\tiny{[P1$\rangle$GM]}}
                \texttt{OTÁZKA: Koľko celkovo vyplnených buniek (X) je vo vašej mriežke? V mojej mriežke je celkovo 16 X. Aký je celkový počet u vás?} \\
            }
        }
    }
    & & \\ \\

    \theutterance \stepcounter{utterance}  
    & & & \multicolumn{2}{p{0.3\linewidth}}{
        \cellcolor[rgb]{0.9,0.9,0.9}{
            \makecell[{{p{\linewidth}}}]{
                \texttt{\tiny{[GM$|$GM]}}
                \texttt{continue} \\
            }
        }
    }
    & & \\ \\

    \theutterance \stepcounter{utterance}  
    & & \multicolumn{4}{p{0.6\linewidth}}{
        \cellcolor[rgb]{0.9,0.9,0.9}{
            \makecell[{{p{\linewidth}}}]{
                \texttt{\tiny{[GM$\rangle$P2]}}
                \texttt{ODPOVEĎ: V mojom piatom riadku je celý riadok vyplnený (X X X X X), teda X vo všetkých piatich stĺpcoch. Vo vašom piatom riadku je vyplnený iba prvý a posledný stĺpec (X ▢ ▢ ▢ X), takže sa líšia. Na základe všetkých doterajších rozdielov je jednoznačné, že naše mriežky NIE sú rovnaké.} \\
\texttt{OTÁZKA: Koľko celkovo vyplnených buniek (X) je vo vašej mriežke? V mojej mriežke je celkovo 16 X. Aký je celkový počet u vás?Svoju odpoveď začnite s "ODPOVEĎ:" a nepridávajte nič iné.} \\
            }
        }
    }
    & \\ \\

    \theutterance \stepcounter{utterance}  
    & & & \multicolumn{4}{p{0.6\linewidth}}{
        \cellcolor[rgb]{0.9,0.9,0.9}{
            \makecell[{{p{\linewidth}}}]{
                \texttt{\tiny{[GM$\langle$P2]}}
                \texttt{ODPOVEĎ: V mojej mriežke je celkovo 13 X (riadky majú 2, 3, 3, 3 a 2 vyplnených buniek). Vo vašej mriežke je 16 X, takže aj celkový počet sa líši. Naše mriežky NIE sú rovnaké.} \\
            }
        }
    }
     \\ \\

    \theutterance \stepcounter{utterance}  
    & & & \multicolumn{2}{p{0.3\linewidth}}{
        \cellcolor[rgb]{0.9,0.9,0.9}{
            \makecell[{{p{\linewidth}}}]{
                \texttt{\tiny{[GM$|$GM]}}
                \texttt{continue} \\
            }
        }
    }
    & & \\ \\

    \theutterance \stepcounter{utterance}  
    & & \multicolumn{4}{p{0.6\linewidth}}{
        \cellcolor[rgb]{0.9,0.9,0.9}{
            \makecell[{{p{\linewidth}}}]{
                \texttt{\tiny{[GM$\rangle$P2]}}
                \texttt{Teraz sa rozhodnite. Čo si myslíte: je vaša mriežka a druhá opísaná mriežka rovnaká? Napíšte "ROZHODNUTIE: rovnaká mriežka" ak si myslíte, že ide o rovnakú mriežku, alebo "ROZHODNUTIE: rôzne mriežky" ak si myslíte, že ide o rôzne mriežky. Nepridávajte nič iné.} \\
            }
        }
    }
    & \\ \\

    \theutterance \stepcounter{utterance}  
    & & & \multicolumn{4}{p{0.6\linewidth}}{
        \cellcolor[rgb]{0.9,0.9,0.9}{
            \makecell[{{p{\linewidth}}}]{
                \texttt{\tiny{[GM$\langle$P2]}}
                \texttt{ROZHODNUTIE: rôzne mriežky} \\
            }
        }
    }
     \\ \\

    \theutterance \stepcounter{utterance}  
    & & & \multicolumn{2}{p{0.3\linewidth}}{
        \cellcolor[rgb]{0.9,0.9,0.9}{
            \makecell[{{p{\linewidth}}}]{
                \texttt{\tiny{[GM$|$GM]}}
                \texttt{continue} \\
            }
        }
    }
    & & \\ \\

    \theutterance \stepcounter{utterance}  
    & & & \multicolumn{2}{p{0.3\linewidth}}{
        \cellcolor[rgb]{0.9,0.9,0.9}{
            \makecell[{{p{\linewidth}}}]{
                \texttt{\tiny{[GM$|$GM]}}
                \texttt{success} \\
            }
        }
    }
    & & \\ \\

    \theutterance \stepcounter{utterance}  
    & & \multicolumn{4}{p{0.6\linewidth}}{
        \cellcolor[rgb]{0.9,0.9,0.9}{
            \makecell[{{p{\linewidth}}}]{
                \texttt{\tiny{[P1$\langle$GM]}}
                \texttt{ODPOVEĎ: V mojej mriežke je celkovo 13 X (riadky majú 2, 3, 3, 3 a 2 vyplnených buniek). Vo vašej mriežke je 16 X, takže aj celkový počet sa líši. Naše mriežky NIE sú rovnaké.} \\
\texttt{Teraz sa rozhodnite. Čo si myslíte: je vaša mriežka a druhá opísaná mriežka rovnaká? Napíšte "ROZHODNUTIE: rovnaká mriežka" ak si myslíte, že ide o rovnakú mriežku, alebo "ROZHODNUTIE: rôzne mriežky" ak si myslíte, že ide o rôzne mriežky. Nepridávajte nič iné.} \\
            }
        }
    }
    & \\ \\

    \theutterance \stepcounter{utterance}  
    & \multicolumn{4}{p{0.6\linewidth}}{
        \cellcolor[rgb]{0.9,0.9,0.9}{
            \makecell[{{p{\linewidth}}}]{
                \texttt{\tiny{[P1$\rangle$GM]}}
                \texttt{ROZHODNUTIE: rôzne mriežky} \\
            }
        }
    }
    & & \\ \\

    \theutterance \stepcounter{utterance}  
    & & & \multicolumn{2}{p{0.3\linewidth}}{
        \cellcolor[rgb]{0.9,0.9,0.9}{
            \makecell[{{p{\linewidth}}}]{
                \texttt{\tiny{[GM$|$GM]}}
                \texttt{continue} \\
            }
        }
    }
    & & \\ \\

    \theutterance \stepcounter{utterance}  
    & & & \multicolumn{2}{p{0.3\linewidth}}{
        \cellcolor[rgb]{0.9,0.9,0.9}{
            \makecell[{{p{\linewidth}}}]{
                \texttt{\tiny{[GM$|$GM]}}
                \texttt{success} \\
            }
        }
    }
    & & \\ \\

\end{supertabular}
}

\end{document}
